% =====================================================================
% TEKNOFEST 2026 - 5G & Yapay Zeka ile Akilli Yol Guvenligi
% FINAL TASARIM RAPORU (FTR)
% Overleaf'te derle: Menu -> Compiler -> XeLaTeX (Turkce karakterler icin)
% =====================================================================
\documentclass[11pt,a4paper]{article}

% --- Turkce + font (XeLaTeX) ---
\usepackage{fontspec}
\usepackage{polyglossia}
\setdefaultlanguage{turkish}
\setmainfont{Latin Modern Roman}

% --- Sayfa, tablo, gorsel, link ---
\usepackage[margin=2.3cm]{geometry}
\usepackage{graphicx}
\usepackage{booktabs}
\usepackage{tabularx}
\usepackage{array}
\usepackage{multirow}
\usepackage{xcolor}
\usepackage{amsmath}
\usepackage{enumitem}
\usepackage[colorlinks=true,linkcolor=blue,urlcolor=blue,citecolor=blue]{hyperref}
\usepackage{caption}
\captionsetup{font=small,labelfont=bf}

% --- Kod blogu icin ---
\usepackage{listings}
\definecolor{kodbg}{RGB}{245,245,245}
\lstset{basicstyle=\ttfamily\footnotesize,backgroundcolor=\color{kodbg},
  breaklines=true,frame=single,numbers=none,columns=fullflexible}

% --- Placeholder makrosu (final sayilar gelince doldur) ---
\newcommand{\PH}[1]{{\color{red}\textbf{[#1]}}}

\title{\textbf{5G \& Yapay Zeka ile Akilli Yol Guvenligi\\ Final Tasarim Raporu}}
\author{Takim Adi: \PH{takim adi} \\ Takim ID: \PH{id} \\ Basvuru ID: \PH{basvuru id}}
\date{Haziran 2026}

\begin{document}
\maketitle
\tableofcontents
\newpage

% =====================================================================
\section{Proje Ozeti (5 Puan)}
% =====================================================================
Bu proje kapsaminda, yol kenari (MOBESE) ve arac-ici (in-cabin) kamera
akislarindan yol guvenligini tehdit eden durumlari gercek zamanli tespit eden,
\textbf{moduler ve cevrim-disi (offline) calisan} bir bilgisayarli goru sistemi
gelistirilmistir. Sistem; arac tipi/renk/plaka tanima, surucu davranis analizi
(telefon kullanma, su icme, esneme, yana/arkaya bakma, sigara icme, emniyet
kemeri ihlali), yolcu konum tespiti, slalom (duzensiz surus) ve acik-kelime
nesne tespiti (\texttt{teknocan}) olmak uzere \PH{N} farkli etiketi uretmektedir.

Cozum, yarismanin Tesla T4 GPU, $<$8\,GB imaj ve 10\,dakika sure kisitlarina
uygun olacak sekilde tasarlanmis; tum cikarim \texttt{results.json} semasina
birebir uyumlu uretilmektedir. Sistemin ayirt edici yani, \textbf{asiri-ogrenme
(overfitting) ve veri sizinti (data leakage) farkindaligidir}: modeller dengeli
veri ve dogru dogrulama protokolleriyle egitilmis, gercek videolarda
\textbf{yanlis-pozitif kontrolu} yapilarak guvenilirligi kanitlanmistir.

% =====================================================================
\section{Veri Seti Olusturulmasi (20 Puan)}
% =====================================================================
Yarisma komitesi ornek video disinda egitim verisi saglamadigindan, modeller
\textbf{acik kaynakli, ticari kullanima uygun lisansli} veri setleriyle
egitilmistir (lisanslar Kaynakca'da belirtilmistir).

\subsection*{Kullanilan Veri Setleri}
\begin{table}[h]
\centering
\small
\begin{tabularx}{\textwidth}{l l c l X}
\toprule
\textbf{Gorev} & \textbf{Veri Seti} & \textbf{Gorsel} & \textbf{Lisans} & \textbf{Kullanim} \\
\midrule
Surucu davranisi & Kaggle DMS (habbas11) & 5957 & Apache 2.0 & sigara/kemer/telefon/goz \\
Arac rengi & VCoR & $\sim$10500 & CC BY 4.0 & 15 renk siniflandirma \\
Plaka & License Plate Recognition & 98798 & CC BY 4.0 & plaka tespiti \\
Negatif (denge) & State Farm c0 (safe) & 2489 & yarisma & arka-plan/negatif \\
Nesne omurgasi & COCO (YOLOv8m) & - & - & telefon/sise/laptop/insan \\
\bottomrule
\end{tabularx}
\caption{Egitimde kullanilan acik kaynak veri setleri.}
\end{table}

\subsection*{Veri Dengeleme (Data Balancing) -- Kritik Tasarim Karari}
On calismalarimizda, yalnizca pozitif ornek iceren (negatif/\,``guvenli surus''
ornegi olmayan) veri setleriyle egitilen modellerin \textbf{``arac-ici sahne =
ihlal'' kisayolunu ogrenerek} temiz videolarda dahi yuksek guvenle yanlis-pozitif
urettigi tespit edilmistir (State Farm 2016 yarismasinda belgelenmis bilinen bir
tuzak). Bu sorunu cozmek icin:
\begin{itemize}[nosep]
  \item Egitim setine \textbf{State Farm c0 (guvenli surus, 2489 gorsel)} ve
        temiz surucu video kareleri \textbf{bos-etiket (background)} olarak
        eklenmis; toplam egitim verisinin \textbf{\%31'i negatif} yapilmistir.
  \item Boylece model ``el agizda = sigara'' gibi kisayollar yerine gercek
        nesneyi ogrenmeye zorlanmistir.
\end{itemize}

\subsection*{Veri Artirma (Data Augmentation)}
Model gurbuzlugu icin domain-randomization tabanli artirma uygulanmistir:
HSV/parlaklik/kontrast jitter (gece/gunduz/IR), Mosaic, MixUp, rastgele
donme/oteleme/olcekleme, rastgele silme (occlusion). Bu, sartnamedeki farkli
FPS, cozunurluk ve isik kosullarina dayaniklilik gereksinimini hedefler.

\subsection*{Egitim/Dogrulama/Test Ayrimi}
\PH{train/val/test oranlari ve gerekce}. Asiri-ogrenmeyi onlemek icin
dogrulama, egitimde gorulmeyen ornekler uzerinde yapilmis; nihai dogrulama
\textbf{egitimde hic kullanilmayan gercek videolarla} gerceklestirilmistir.

% =====================================================================
\section{Yapay Zeka Cozumu (50 Puan)}
% =====================================================================

\subsection{Problemin Analizi (15 Puan)}
Video uzerinden tespit yaparken karsilasilan temel problemler ve cozum
yaklasimimiz:
\begin{itemize}[nosep]
  \item \textbf{Asiri-ogrenme / veri sizintisi:} Surucu veri setleri video
        kareleridir; komsu kareler train/val'e dagilirsa dogrulama skoru
        yaniltici sekilde sisar (\%99 gorunup gercekte \%38 olabilir). Cozum:
        dengeli veri + gercek-video dogrulamasi.
  \item \textbf{Kamera acisi belirsizligi:} DMS kamerasi onden veya yan (profil)
        olabilir. Basit landmark-geometri yaw tahmini profilde basarisiz oldugundan
        \textbf{appearance-tabanli 6DRepNet} (3B kafa pozu) + per-video referans
        kalibrasyonuna gecilmistir.
  \item \textbf{Isik degisimi, hareket bulanikligi, okluzyon:} domain-randomization
        artirma + zamansal oylama (track-level voting) ile ele alinmistir.
  \item \textbf{Yanlis-pozitif:} Zayif/ezberci modeller devre disi birakilmis,
        guvenilir kaynaklara (dengeli COCO) oncelik verilmistir.
\end{itemize}

\subsection{Cozum Mimarisi (15 Puan)}
Sistem, ham videodan \texttt{results.json} ciktisina kadar moduler bir hatta
(pipeline) calisir. Her katman \textbf{izole}dir: biri hata verse dahi ana akis
gecerli cikti uretir (sartnamedeki hata-yonetimi gereksinimi).

\begin{figure}[h]
\centering
\fbox{\parbox{0.92\textwidth}{\centering \PH{Mimari diyagram buraya} \\[4pt]
\footnotesize Video $\rightarrow$ kare ornekleme ($\sim$5 FPS) $\rightarrow$
YOLOv8m (COCO) + ByteTrack izleme $\rightarrow$ track-level cogunluk oylama
$\rightarrow$ \{Arac: tip/renk(VCoR)/plaka(YOLO+EasyOCR) $|$ Davranis:
COCO telefon-su, MediaPipe esneme/yorgunluk, 6DRepNet bakma, sigara/kemer modeli
$|$ Yolcu $|$ Slalom $|$ teknocan (YOLO-World)\} $\rightarrow$ tespit dedup
$\rightarrow$ \texttt{results.json}}}
\caption{Sistemin kusbakisi mimarisi (placeholder -- net diyagram eklenecek).}
\end{figure}

\subsection{Cozum Detaylari (20 Puan)}
\textbf{Omurga ve izleme:} YOLOv8m (COCO on-egitimli) nesne tespiti +
ByteTrack ile coklu-nesne takibi. Track boyunca \textbf{zamansal cogunluk
oylamasi} tek-kare yanlis-pozitifleri eler, kacan kareleri koprular
(precision \& recall artisi).

\textbf{Arac analizi:}
\begin{itemize}[nosep]
  \item Tip: bbox oran sezgisi + COCO sinifi.
  \item Renk: VCoR uzerinde egitilmis YOLOv8m-cls (15 renk $\rightarrow$ 9 etiket).
  \item Plaka: License Plate (mAP \PH{X}) ile bbox $\rightarrow$ 4-nokta perspektif
        duzeltme $\rightarrow$ EasyOCR (\texttt{tr}, allowlist) $\rightarrow$ TR plaka
        regex dogrulamasi + karakter-bazli oylama.
\end{itemize}

\textbf{Surucu davranisi:}
\begin{itemize}[nosep]
  \item Telefon/su/bilgisayar: COCO (dengeli, milyonlarca ornek) -- en guvenilir
        kaynak.
  \item Esneme/yorgunluk: MediaPipe Face Mesh -- MAR (agiz) ve EAR (goz) sezgileri,
        egitimsiz, cevrim-disi.
  \item Yana/arkaya bakma: \textbf{6DRepNet} 3B kafa pozu + ardisik-sureklilik
        sarti + per-video referans (kamera acisina dayanikli).
  \item Sigara/emniyet kemeri: Kaggle DMS (Apache 2.0) uzerinde negatif-dengeli
        fine-tune YOLOv8m.
\end{itemize}

\textbf{Slalom:} ByteTrack yorungesinden yanal salinim analizi (egitimsiz).\\
\textbf{teknocan:} YOLO-World acik-kelime tespiti (gomulu embedding, cevrim-disi).

\textbf{Donanim/Yazilim:} PyTorch + Ultralytics YOLOv8, OpenCV, MediaPipe,
EasyOCR, 6DRepNet; Docker (\texttt{nvidia/cuda:12.1.0-base-ubuntu22.04}),
Tesla T4. Tum agirliklar imaja gomulu -- \textbf{calisma aninda internet kapali}.

% =====================================================================
\section{Cozumun Sinanmasi (20 Puan)}
% =====================================================================
Modeller hem standart metriklerle hem de \textbf{egitimde gorulmeyen gercek
videolarla} sinanmistir. ``Cozumumuze neden guveniyoruz?'' sorusunun yaniti
asagidaki kanitlardir.

\subsection*{Model Metrikleri (Dogrulama)}
\begin{table}[h]
\centering
\small
\begin{tabularx}{\textwidth}{l c c c c X}
\toprule
\textbf{Model} & \textbf{Precision} & \textbf{Recall} & \textbf{mAP50} & \textbf{mAP50-95} & \textbf{Not} \\
\midrule
sigara+kemer (DMS) & \PH{P} & \PH{R} & \PH{mAP} & \PH{mAP95} & negatif-dengeli, durust \\
Plaka & 0.986 & 0.955 & 0.972 & 0.678 & 98K gorsel \\
Renk (VCoR) & - & - & - & - & top1 acc \%87.7 \\
\bottomrule
\end{tabularx}
\caption{Egitilen modellerin dogrulama metrikleri.}
\end{table}

\subsection*{Cevrim-Disi (Offline) Docker Testi}
Sistem \texttt{docker run --network none} ile (internet tamamen kapali, T4)
gercek 4K trafik videosunda calistirilmis; gecerli \texttt{results.json}
uretilmis, cokme yasanmamistir. Imaj boyutu \textbf{4.5\,GB} ($<$8\,GB limit),
isleme suresi T4'te 10\,dk limitinin cok altindadir.

\subsection*{Gercek Video Testi -- Yanlis-Pozitif Kanitlari}
En kritik dogrulama: temiz ve ihlalli videolarin ayirt edilmesi.
\begin{table}[h]
\centering
\small
\begin{tabularx}{\textwidth}{l X l}
\toprule
\textbf{Test Videosu} & \textbf{Beklenen} & \textbf{Sonuc} \\
\midrule
goodmax (temiz surucu) & ihlal YOK & \PH{temiz mi?} \\
badmax (telefon/sigara/bakma) & ihlaller & \PH{yakalandi mi?} \\
Gercek plaka videolari & plaka okuma & 03ACU808, 64EP367 (basarili) \\
\bottomrule
\end{tabularx}
\caption{Gercek videolarla yanlis-pozitif ve dogru-pozitif sinamasi.}
\end{table}

\subsection*{Muhendislik Kalitesi}
\PH{N} birim testi (TDD), izole hata-yonetimli katmanlar, JSON sema dogrulayici,
ASCII-safe etiket garantisi, ortam-tespiti iceren manipulasyon \textbf{yok}
(anti-cheat uyumlu).

% =====================================================================
\section{Kaynakca (5 Puan)}
% =====================================================================
\begin{thebibliography}{9}
\footnotesize
\bibitem{dms} Kaggle DMS Dataset (habbas11), Apache 2.0.
  \url{https://www.kaggle.com/datasets/habbas11/dms-driver-monitoring-system}
\bibitem{vcor} VCoR Vehicle Color Recognition Dataset, Kezebou et al., MDPI AI 2021, CC BY 4.0.
\bibitem{plate} License Plate Recognition Dataset (Roboflow Universe), CC BY 4.0.
\bibitem{statefarm} State Farm Distracted Driver Detection, Kaggle (yarisma).
\bibitem{yolov8} Jocher et al., Ultralytics YOLOv8, 2023.
\bibitem{bytetrack} Zhang et al., ByteTrack: Multi-Object Tracking by Associating Every Detection Box, ECCV 2022.
\bibitem{6drepnet} Hempel et al., 6D Rotation Representation for Unconstrained Head Pose Estimation, ICIP 2022.
\bibitem{mediapipe} Lugaresi et al., MediaPipe: A Framework for Building Perception Pipelines, 2019.
\bibitem{yoloworld} Cheng et al., YOLO-World: Real-Time Open-Vocabulary Object Detection, CVPR 2024.
\bibitem{easyocr} JaidedAI, EasyOCR, Apache 2.0.
\end{thebibliography}

\end{document}
